%% ═══════════════════════════════════════════════════════════════════
\chapter*{Conclusioni}
\addcontentsline{toc}{chapter}{Conclusioni}
%% ═══════════════════════════════════════════════════════════════════

%% ─── PROMPT LLM ─────────────────────────────────────────────────
%% Scrivi le conclusioni della tesi. Tono: critico, tagliente,
%% intellettualmente onesto. Non è un riassunto: è una riflessione
%% finale che sfida il lettore e lascia un'impressione duratura.
%% Struttura in 5 blocchi distinti:
%%
%% BLOCCO 1 -- RISPOSTA ALLE TRE RQ (sintetica, 1 paragrafo ciascuna):
%%   RQ1: Il sistema NDC è finanziariamente sostenibile MA non garantisce
%%        adeguatezza alle generazioni future con carriere discontinue
%%        e crescita economica stagnante. Il trade-off è reale.
%%   RQ2: I correttivi tecnici esistono (aggiornamento automatico
%%        coefficienti, pensione di garanzia, separazione
%%        previdenza/assistenza) ma sono bloccati dall'economia politica,
%%        non dalla tecnica. La robot tax è la domanda aperta del futuro.
%%   RQ3: Il singolo lavoratore informato ha strumenti concreti:
%%        fondo pensione negoziale, TFR investito, PAC su ETF,
%%        investimento in capitale umano. Ma non può essere questa
%%        l'unica risposta sistemica.
%%
%% BLOCCO 2 -- LA CONTRADDIZIONE CENTRALE (1 pagina, tono critico):
%%   "Un sistema progettato per garantire un tenore di vita dignitoso
%%    in vecchiaia potrebbe non farlo per milioni di lavoratori Gen Z.
%%    Questa non è una profezia catastrofista: è l'algebra del NDC
%%    applicata a 30 anni di crescita italiana quasi nulla."
%%   La sostenibilità finanziaria è necessaria ma non sufficiente.
%%   I conti tornano. Le pensioni, per molti, potrebbero non bastare.
%%
%% BLOCCO 3 -- IL PARADOSSO DEMOCRATICO (provocatorio):
%%   "In Italia, le riforme pensionistiche razionali vengono bloccate
%%    dall'elettore mediano di 52 anni. I giovani, i più penalizzati,
%%    votano meno di tutti. Il sistema pensionistico italiano non è
%%    fallito per ragioni tecniche: è stato catturato politicamente
%%    da chi dal sistema ha già tratto il massimo vantaggio."
%%   Cita: Fornero (2018), dati astensionismo giovanile.
%%
%% BLOCCO 4 -- LA DOMANDA APERTA SULL'AUTOMAZIONE:
%%   "Se nei prossimi 20 anni l'AI sostituirà una quota significativa
%%    del lavoro umano, il PAYG entra in crisi strutturale indipendentemente
%%    dall'efficienza del NDC. Non è una questione tecnica-previdenziale:
%%    è una questione filosofica su chi ha diritto alla ricchezza prodotta
%%    dalle macchine. Il Contributo Automazione è una risposta possibile.
%%    Ma la domanda vera è più profonda: il contratto sociale del welfare
%%    è stato scritto nel Novecento, per un mondo di lavoratori dipendenti
%%    a tempo pieno. Quel mondo non esiste più."
%%
%% BLOCCO 5 -- LA DOMANDA NORMATIVA FINALE (la più tagliente):
%%   "È giusto che la responsabilità di una pensione adeguata ricada
%%    sul singolo? In un sistema liberale coerente, la risposta è sì:
%%    lo Stato garantisce la soglia minima, il resto è responsabilità
%%    individuale. Ma questa risposta è accettabile solo se il lavoratore
%%    ha avuto accesso a informazioni complete, strumenti accessibili
%%    e una crescita economica che rendesse i contributi produttivi.
%%    In Italia, nessuna delle tre condizioni è pienamente soddisfatta.
%%    Quindi la domanda rimane aperta, e con essa la tensione
%%    irrisolta tra rigore contabile e giustizia sociale."
%% ────────────────────────────────────────────────────────────────

%% ─── BLOCCO 4 — LA DOMANDA APERTA SULL'AUTOMAZIONE ─────────────────
\section*{La domanda che il sistema non ha ancora fatto}
\addcontentsline{toc}{section}{La domanda che il sistema non ha ancora fatto}

\noindent
Il quarto capitolo ha documentato il paradosso strutturale: un'economia
in cui la ricchezza cresce per effetto dell'automazione vede simultaneamente
erodersi il proprio monte contributivo, perché i profitti da capitale non
entrano nel circuito PAYG.
Non è una proiezione astratta: il FMI stima che il 40\% dei lavori nei
paesi avanzati sia già esposto all'automazione cognitiva \cite{IMF2024},
e Acemoglu e Restrepo documentano che ogni robot installato nei settori
manifatturieri comprime il tasso di occupazione locale di circa 0,34
punti percentuali \cite{AcemogluRestrepo2020}.
Il denominatore della formula NDC, il monte contributivo aggregato, non
è solo una variabile demografica: è anche una variabile tecnologica.

Il Contributo Automazione è la risposta parziale che il dibattito italiano
ha prodotto: tassare i sistemi che sostituiscono il lavoro per finanziare
il welfare che quel lavoro avrebbe sostenuto \cite{Bacchiocchi2025}.
La logica è corretta. Ma funziona come risposta alla crisi fiscale del
breve periodo, non come risposta alla domanda più profonda.

Quella domanda è filosofica.
Il contratto sociale su cui poggia ogni sistema a ripartizione è stato
scritto nel secondo dopoguerra per un mondo di lavoratori dipendenti a
tempo pieno, con carriere lineari, salari in crescita strutturale e un
orizzonte demografico che prometteva più lavoratori domani di quanti
ce ne fossero oggi.
Nessuna di queste premesse regge nell'Italia del 2024: le carriere sono
discontinue per struttura, i salari reali stagnano da un quarto di secolo,
e l'indice di dipendenza demografica supererà il 57\% entro il 2040
\cite{MEFRGS2025, Fornero2018}.

Se nei prossimi vent'anni l'AI cognitiva automaterà il 25--40\% delle
attività a media qualificazione, il PAYG entrerà in crisi strutturale
indipendentemente dall'efficienza della formula NDC: sarà un sistema
progettato per un'economia del lavoro che opera in un'economia del capitale
\cite{AcemogluJohnson2023}.
La risposta strutturalmente coerente è quella che Bosi aveva già
delineato con la separazione tra funzione previdenziale e funzione
assistenziale \cite{Bosi2018}: finanziare con la fiscalità generale la
componente assistenziale del welfare, quella che garantisce la soglia
minima a prescindere dalla storia contributiva; lasciare alla
contribuzione sul reddito da lavoro la componente proporzionale.
Un sistema con questa architettura sopravviverebbe all'automazione
perché la sua base di finanziamento non è legata esclusivamente al
reddito da lavoro.

Questa non è una riforma disponibile nel breve periodo: richiede un
accordo politico sulla struttura del welfare che in Italia non ha mai
trovato la maggioranza necessaria.
È però la direzione verso cui le prossime riforme dovrebbero guardare,
se avranno il coraggio di fare la domanda giusta invece di continuare
a rispondere a quella già superata.


%% ─── BLOCCO 5 — LA DOMANDA NORMATIVA FINALE ──────────────────────────
\section*{La responsabilità individuale come principio e come alibi}
\addcontentsline{toc}{section}{La responsabilità individuale come principio e come alibi}

\noindent
C'è una domanda normativa che questa tesi ha sollevato in ogni capitolo
senza mai chiuderla: è giusto che la responsabilità di una pensione
adeguata ricada sul singolo lavoratore?

In un framework liberale coerente, la risposta è sì.
Lo Stato garantisce la soglia minima attraverso la pensione contributiva
di garanzia e l'assegno sociale.
Tutto ciò che supera quella soglia è responsabilità individuale: chi ha
lavorato, versato contributi, aderito a un fondo pensione e gestito
il TFR in modo razionale ottiene una vecchiaia economicamente sostenibile
perché ha esercitato la propria agentività con consapevolezza.
La neutralità attuariale è, in questa lettura, un principio di giustizia
distributiva prima ancora che un vincolo tecnico: a parità di contributi,
a parità di pensione \cite{Bosi2018}.

Questa risposta è coerente. Ma è accettabile solo se tre condizioni sono
soddisfatte simultaneamente.

La prima condizione è l'informazione.
Il lavoratore deve sapere, con sufficiente anticipo, quale pensione si
aspetta, a quale tasso di sostituzione si troverà e quanto deve integrare
per colmare il divario.
La busta arancione, lo strumento di comunicazione previdenziale
individuale che la riforma Dini aveva previsto, non è mai stata
implementata in modo sistematico; la piattaforma INPS fornisce proiezioni
di difficile interpretazione per chi non ha formazione attuariale; le
indagini di settore mostrano che la maggioranza dei lavoratori dipendenti
italiani non conosce il proprio tasso di sostituzione atteso
\cite{Fornero2018, Lavoce2024a}.
Chiedere scelte informate in assenza di informazione non è un principio
liberale: è una delega in bianco.

La seconda condizione è l'accessibilità degli strumenti.
La deduzione fiscale sui contributi al fondo pensione, fino a 5.164,57
euro annui, è uno strumento potente: trasforma il risparmio complementare
da sacrificio aggiuntivo in redistribuzione fiscale già incorporata nel
sistema.
Ma il mercato dei fondi pensione italiani non è uniforme: i piani
individuali pensionistici distribuiti dalle reti bancarie e assicurative
presentano in molti casi costi di gestione tra i più alti in Europa,
asset allocation troppo conservativa per orizzonti trentennali e una
governance che difende i margini dei distributori più degli interessi
degli aderenti \cite{Bosi2018, Lavoce2024a}.
Gli strumenti esistono. L'accesso equo a strumenti di qualità, no.

La terza condizione è la crescita.
La logica attuariale produce risultati giusti se il sistema ha remunerato
adeguatamente i contributi versati nel corso della vita lavorativa.
Il tasso di capitalizzazione nozionale del sistema NDC è la media mobile
quinquennale del PIL nominale: con una crescita reale italiana attestata
attorno allo 0\% tra il 2000 e il 2024 \cite{MEFRGS2025}, il rendimento
implicito del sistema è rimasto sistematicamente al di sotto di qualsiasi
alternativa di risparmio ragionevolmente accessibile.
Non si può chiedere responsabilità individuale per un risultato che la
macroeconomia del paese non ha mai reso producibile.

In Italia, nessuna delle tre condizioni è stata garantita in modo
sistematico.

Questo non significa abbandonare il principio di responsabilità
individuale: significa chiedere allo Stato di fare quello che, in un
framework liberale, è suo compito fare prima che i giocatori entrino in
campo, garantire che le regole siano uguali, chiare e accessibili a tutti.
Informazione pensionistica trasparente e obbligatoria, riduzione dei costi
nei fondi negoziali attraverso la concorrenza e la regolazione, separazione
netta tra funzione previdenziale e funzione assistenziale: nessuna di
queste misure richiede di abbandonare la neutralità attuariale.
Richiedono che quella neutralità valga davvero per tutti, e non soltanto
per chi aveva già le condizioni per sfruttarla \cite{Fornero2018, Artoni2014}.

La tensione tra rigore contabile e giustizia sociale che questa tesi ha
esaminato non si risolve in queste pagine.
Il sistema NDC è tecnicamente corretto.
Le proiezioni RGS sono credibili.
E al tempo stesso un sistema che scarica sull'individuo il rischio
demografico e di longevità, senza garantirgli informazione, strumenti
e crescita per gestirlo, non è un sistema equo: è un sistema che ha
razionalizzato la propria inadeguatezza in linguaggio tecnico.

Non è una contraddizione da sciogliere qui.
È il problema che la generazione che entra oggi nel mercato del lavoro
eredita, insieme ai contributi da versare.
